% Viva presentation — EEEM004, Fiyinfoluwa Akano, September 2026.
% 12 slides (handbook maximum), sized for a 15-20 minute talk.
% Compile: xelatex viva_deck.tex   (run twice)
\documentclass[aspectratio=169,11pt]{beamer}

\usetheme{default}
\usecolortheme{seahorse}
\setbeamertemplate{navigation symbols}{}
\setbeamertemplate{footline}{%
  \hfill\scriptsize\insertframenumber/\inserttotalframenumber\hspace{2mm}\vspace{2mm}}
\setbeamertemplate{itemize items}{\textbullet}
\setbeamerfont{frametitle}{series=\bfseries,size=\large}
\setbeamercolor{frametitle}{fg=black,bg=white}
\setbeamercolor{normal text}{fg=black}
\setbeamercolor{structure}{fg=black!70}

\usepackage{fontspec}
\IfFontExistsTF{Times New Roman}{\setmainfont{Times New Roman}}{}
\usefonttheme{serif}
\usepackage{booktabs}
\usepackage{graphicx}
\graphicspath{{../dissertation/}}

\title{Evaluating Reinforcement Learning Algorithms\\for Portfolio Trading}
\subtitle{A Reinforcement Learning Study on the S\&P 500 ETF (SPY)}
\author{Fiyinfoluwa Akano \\ \small URN 6962514 · MSc Artificial Intelligence}
\institute{Supervisor: Dr Cuong Nguyen \\ University of Surrey}
\date{September 2026}

\begin{document}

% ---------- 1. Title ----------
\begin{frame}
  \titlepage
  \begin{center}
    \includegraphics[width=0.22\textwidth]{figs/surrey_logo.png}
  \end{center}
\end{frame}

% ---------- 2. Problem ----------
\begin{frame}{The problem}
  An investor holding a stock-index fund faces the same decision each
  day: buy more, sell part of the position, or do nothing.
  \vspace{4pt}
  \begin{itemize}
    \item \textbf{Question 1 — performance.} Can reinforcement learning
      make these decisions better than buying and holding, after
      transaction costs?
    \item \textbf{Question 2 — behaviour.} Does the learned agent
      actually use its market information? A market that mostly rises
      rewards constant buying, so a profitable agent has not
      necessarily learned anything.
    \item \textbf{Question 3 — data.} Is the historical record of one
      asset large enough to train these networks at all?
  \end{itemize}
  \vspace{4pt}
  Three algorithms compared: REINFORCE and Deep Q-learning (DQN),
  implemented from scratch, and PPO from a library. The dissertation
  answers with \textbf{two studies}: real data first, then corrected
  data.
\end{frame}

% ---------- 3. Environment ----------
\begin{frame}{The trading environment}
  \begin{columns}[T]
    \begin{column}{0.55\textwidth}
      \begin{itemize}
        \item Monthly episodes of daily SPY bars (18--23 trading days),
          \$10,000 starting capital, any open position sold on the
          final day.
        \item \textbf{State}: nine features --- five market (return,
          momentum, volatility, trend gap, range) and four portfolio
          (cash, exposure, unrealised PnL, episode clock).
        \item \textbf{Actions}: hold, buy, or sell a fixed
          \(0.25\,C_0\); invalid actions masked.
        \item \textbf{Reward}: change in wealth after a 0.05\% fee.
          Risk-aware variant subtracts \(\lambda C_0\,\Delta dd_t\).
      \end{itemize}
    \end{column}
    \begin{column}{0.42\textwidth}
      \includegraphics[width=\textwidth]{media/media/image1_ch3.png}
    \end{column}
  \end{columns}
\end{frame}

% ---------- 4. Methodology ----------
\begin{frame}{How the evaluation earns trust}
  \begin{itemize}
    \item \textbf{Chronological splits.} Train 2018--2022, validate
      2023, test 2024--2025. No shuffling, so no look-ahead leakage.
    \item \textbf{Validation-only decisions.} Trade size, risk-penalty
      weight, and fee checks were all selected on validation data.
      Every test set was touched once.
    \item \textbf{Six seeds per configuration}, spread reported next to
      every mean. One good run proves nothing.
    \item \textbf{21 verification tests.} They caught a benchmark bug
      worth \$38.95 per month --- enough to flip the headline
      conclusion.
    \item \textbf{State-dependence test.} A policy that picks different
      actions at steps with the \emph{same action mask} must be reading
      its state. Reported per seed, next to the returns.
  \end{itemize}
  \vspace{2pt}
  The last test separates a policy that responds to the market from one
  that follows a fixed pattern --- the returns alone cannot.
\end{frame}

% ---------- 5. First study ----------
\begin{frame}{First study: real data alone --- a negative result}
  Trained on the 60 real months of 2018--2022, tested on 2024--2025.
  \vspace{4pt}
  \begin{columns}[T]
    \begin{column}{0.52\textwidth}
      \small
      \begin{tabular}{@{}lrr@{}}
        \toprule
        Policy & \(\Delta W\) (\$/mo) & State-dep. \\
        \midrule
        Buy-and-hold & +180.25 & --- \\
        Always-buy   & +179.75 & --- \\
        REINFORCE    & +171.51 & 0/6 \\
        PPO          & +127.35 & 0/6 \\
        Deep Q-learning & +74.68 & \textbf{6/6} \\
        \bottomrule
      \end{tabular}
    \end{column}
    \begin{column}{0.46\textwidth}
      \includegraphics[width=\textwidth]{media4/media/image4.png}
    \end{column}
  \end{columns}
  \vspace{6pt}
  No agent beat buy-and-hold. The methods that earned the most ignored
  their state on every seed; the only method that used its state earned
  the least. The contradiction needed an explanation.
\end{frame}

% ---------- 6. Diagnosis ----------
\begin{frame}{The diagnosis: the data, not the algorithms}
  \begin{itemize}
    \item \textbf{60 episodes against 18,000 parameters.} A network
      that size can memorise every training month; nothing forces it to
      learn a rule that transfers.
    \item \textbf{The months lean upward.} 29 of 60 training months
      rose more than 2\%; only 16 fell. Under that distribution,
      always-buy is a nearly optimal answer that needs no state
      information --- the same failure as a classifier trained on a
      99\%-one-class dataset.
    \item \textbf{Why DQN was worst.} Value estimation needs many
      visits to similar states under a stable distribution. Sixty
      shifting months provide neither, so DQN conditioned on noisy
      estimates.
  \end{itemize}
  \vspace{4pt}
  None of this can be tested on the real data, because the real data is
  the problem. Testing it needs control over volume, balance, and
  conditions --- which is what the main study builds.
\end{frame}

% ---------- 7. The fix ----------
\begin{frame}{The main study: controlled data, right-sized model}
  \begin{columns}[T]
    \begin{column}{0.55\textwidth}
      \begin{itemize}
        \item Market generator with three conditions --- up, down, flat
          --- calibrated on the 60 real training months \emph{only}.
        \item 3,000 balanced training episodes, 300 validation, 600
          test. About 60,000 transitions per pass.
        \item Network cut to 64 and 32 units: \(\approx\)2,800
          parameters. Data now outnumbers parameters \(\sim\)20:1.
        \item On the balanced test set, \textbf{no fixed pattern
          profits}: buy-and-hold loses \$60.29 per episode; never
          trading earns \$0. Any profit requires reading the state.
      \end{itemize}
    \end{column}
    \begin{column}{0.42\textwidth}
      \small
      \begin{tabular}{@{}lrrr@{}}
        \toprule
        & Months & Drift & Vol. \\
        \midrule
        Up   & 29 & +0.244\% & 0.97\% \\
        Down & 16 & −0.358\% & 2.00\% \\
        Flat & 15 & +0.012\% & 1.10\% \\
        \bottomrule
      \end{tabular}
      \vspace{4pt}

      \scriptsize Fitted daily parameters per condition; generated ATR
      matches the real mean within 4\%.
    \end{column}
  \end{columns}
\end{frame}

% ---------- 8. Main result ----------
\begin{frame}{Main result: the ranking inverts}
  \centering
  \includegraphics[width=0.76\textwidth]{figs5/fig5_sim.pdf}
  \vspace{2pt}
  \begin{itemize}\small
    \item \textbf{DQN: +\$126.94 per episode} (seed std \$7.30), all
      six seeds profitable and state-dependent. It keeps 99\% of the
      up-market profit and cuts the down-market loss by 84\%.
    \item REINFORCE and PPO collapse to corner policies --- never trade
      or always buy. One REINFORCE seed in six found the conditioned
      policy (+\$83.13), so it is reachable but not reliably found.
  \end{itemize}
\end{frame}

% ---------- 9. Robustness ----------
\begin{frame}{The result survives its assumptions}
  \textbf{Risk-aware reward} (\(\lambda=0.25\), selected on validation):
  \begin{itemize}
    \item Down-market loss falls 73\% (−\$87.28 to −\$23.26) for a
      1.5\% cost in mean return; Sharpe rises 0.285 to 0.332.
    \item The agent holds less \emph{when the market falls}, not less
      overall --- the selective behaviour the first study could not
      produce.
  \end{itemize}
  \vspace{4pt}
  \textbf{Transaction fees} swept 0 to 50 basis points:
  \begin{itemize}
    \item DQN stays profitable at every level: +\$95.27 per episode at
      50 bps, ten times the study's assumption. Buy-and-hold turns
      negative by 10 bps on balanced data.
    \item Trades fall only from 5.5 to 3.7 per episode ---
      condition-selective trades are worth more than they cost.
  \end{itemize}
  \vspace{4pt}
  \textbf{Drawdown state feature}: +\$4.37 change against a \$7 seed
  std --- information only helps when the objective gives a reason to
  use it.
\end{frame}

% ---------- 10. Transfer ----------
\begin{frame}{Transfer: 180 real months the training never saw}
  \centering
  \includegraphics[width=0.68\textwidth]{figs5/fig5_transfer.pdf}
  \vspace{2pt}
  \begin{itemize}\small
    \item Every real month outside calibration: 2006--2017 and
      2023--2025, including the 2008 crisis. Profitable on all six
      seeds (+\$60.28/mo); \textbf{Sharpe 0.224 vs buy-and-hold's
      0.198} at three-quarters of the exposure.
    \item \textbf{October 2008}: all six agents held zero exposure;
      buy-and-hold lost \$1,665. Crisis total: −\$180 vs −\$3,046.
    \item The honest half: raw return is lower (+\$60 vs +\$80), and
      the agents sat out the Dec 2008 / Mar 2009 rebounds. The claim is
      risk-adjusted, not absolute.
  \end{itemize}
\end{frame}

% ---------- 11. Limitations ----------
\begin{frame}{Limitations and future work}
  \textbf{Limitations}
  \begin{itemize}
    \item One asset; findings do not automatically generalise past
      broad index trading.
    \item The generator is deliberately simple: returns independent
      within a condition, no mid-month regime change. The transfer test
      measures the cost of that gap directly.
    \item Six seeds, no formal significance testing; policy-gradient
      methods were not separately tuned (entropy coefficient 0).
  \end{itemize}
  \vspace{4pt}
  \textbf{Future work}
  \begin{itemize}
    \item Regime switching within episodes (hidden Markov chain) and
      volatility clustering; walk-forward recalibration.
    \item Risk-relative objectives that cannot score well by
      abstention; more seeds with bootstrap confidence intervals;
      other assets.
  \end{itemize}
\end{frame}

% ---------- 12. Conclusions ----------
\begin{frame}{Conclusions}
  \begin{enumerate}
    \item A verified evaluation framework: validation-only decisions,
      six seeds, 21 automated tests.
    \item A behavioural test that separates state-driven policies from
      fixed patterns --- and changed the interpretation of every
      result.
    \item An explained negative result: at real single-asset volumes,
      none of the algorithms learns a genuine policy.
    \item A calibrated simulator that turns market condition into a
      controlled variable.
    \item Evidence that \textbf{the amount and balance of training
      data --- not the choice of algorithm --- set the limit}: fixing
      the data inverted the ranking, and the learned policy carried its
      defensive behaviour into two decades of real markets it had never
      seen.
  \end{enumerate}
  \vspace{6pt}
  \centering
  What these methods learn is decided by the data they are given.
\end{frame}

\end{document}
